\documentclass[12pt,a4paper]{article}

\usepackage[T1]{fontenc}
\usepackage{amsmath}
\usepackage{amssymb}
\usepackage{amsthm}
\usepackage[margin=2.5cm]{geometry}
\usepackage{hyperref}
\usepackage[dvipsnames]{xcolor}
\usepackage{tikz}
\usepackage{pgfplots}
\pgfplotsset{compat=1.18}
\usepackage{tcolorbox}
\usepackage{booktabs}
\usepackage{array}
\usepackage{enumitem}
\usepackage{fancyhdr}
\usepackage{titlesec}
\usepackage{microtype}
\usepackage{setspace}

% Set headheight for fancyhdr
\setlength{\headheight}{14pt}

% TikZ Libraries
\usetikzlibrary{arrows.meta,calc,positioning,shapes.geometric,decorations.pathreplacing,backgrounds,fit,matrix}

% Geometry and Headers
\pagestyle{fancy}
\fancyhf{}
\rhead{\small Linear Algebra: Row Operations}
\lhead{\small Lecture Notes}
\cfoot{\thepage}

% Custom Math Macros
\newcommand{\R}{\mathbb{R}}
\newcommand{\vv}[1]{\mathbf{#1}}
\newcommand{\vzero}{\mathbf{0}}
\newcommand{\Span}{\text{Span}}
\newcommand{\rank}{\text{rank}}
\newcommand{\nullity}{\text{nullity}}
\newcommand{\rref}{\text{rref}}
\newcommand{\adj}{\text{adj}}

% TColorBox Styles
\tcbset{
	defstyle/.style={
		colback=Blue!5,
		colframe=Blue!75!black,
		arc=0mm,
		boxrule=1.5pt,
		left=5pt,
		right=5pt,
		top=5pt,
		bottom=5pt,
		fonttitle=\bfseries,
		title={#1}
	},
	thmstyle/.style={
		colback=Cyan!5,
		colframe=Cyan!75!black,
		arc=0mm,
		boxrule=1.5pt,
		left=5pt,
		right=5pt,
		top=5pt,
		bottom=5pt,
		fonttitle=\bfseries,
		title={#1}
	},
	exstyle/.style={
		colback=Yellow!10,
		colframe=Orange!75!black,
		arc=0mm,
		boxrule=1.5pt,
		left=5pt,
		right=5pt,
		top=5pt,
		bottom=5pt,
		fonttitle=\bfseries,
		title={#1}
	},
	remstyle/.style={
		colback=Red!5,
		colframe=Red!75!black,
		arc=0mm,
		boxrule=1.5pt,
		left=5pt,
		right=5pt,
		top=5pt,
		bottom=5pt,
		fonttitle=\bfseries,
		title={#1}
	}
}

% Environment Definitions
\newtcolorbox{definition}[1]{defstyle={Definition: #1}}
\newtcolorbox{theorem}[1]{thmstyle={Theorem: #1}}
\newtcolorbox{example}[1]{exstyle={Example: #1}}
\newtcolorbox{remark}[1]{remstyle={Remark: #1}}
\newtcolorbox{proposition}[1]{thmstyle={Proposition: #1}}
\newtcolorbox{corollary}[1]{thmstyle={Corollary: #1}}

% Title Formatting
\titleformat{\section}{\Large\bfseries\color{Blue!75!black}}{}{0em}{}[\vspace{0.2cm}\hrule\vspace{0.2cm}]
\titleformat{\subsection}{\large\bfseries\color{Cyan!75!black}}{}{0em}{}

% Line Spacing
\onehalfspacing

% Title Page
\title{\textbf{Linear Algebra: Matrix Operations and Row Reduction}}
\author{Lecture Notes}
\date{\today}

\begin{document}
	
	\maketitle
	
	\tableofcontents
	
	\newpage
	
	\section{Introduction to Matrix Operations}
	
	Matrices are fundamental objects in linear algebra. A matrix is a rectangular array of numbers arranged in rows and columns.
	
	\begin{definition}{Matrix}
		An $m \times n$ matrix is a rectangular array of $mn$ numbers arranged in $m$ rows and $n$ columns:
		\[
		A = \begin{pmatrix}
			a_{11} & a_{12} & \cdots & a_{1n} \\
			a_{21} & a_{22} & \cdots & a_{2n} \\
			\vdots & \vdots & \ddots & \vdots \\
			a_{m1} & a_{m2} & \cdots & a_{mn}
		\end{pmatrix}
		\]
		where $a_{ij} \in \R$ is the entry in row $i$ and column $j$. We denote this matrix as $A \in \R^{m \times n}$.
	\end{definition}
	
	\subsection{Matrix Addition and Scalar Multiplication}
	
	\begin{definition}{Matrix Addition}
		Let $A, B \in \R^{m \times n}$. The sum $A + B$ is the $m \times n$ matrix whose $(i,j)$-entry is $a_{ij} + b_{ij}$.
	\end{definition}
	
	\begin{definition}{Scalar Multiplication}
		Let $A \in \R^{m \times n}$ and $c \in \R$. The scalar multiple $cA$ is the $m \times n$ matrix whose $(i,j)$-entry is $ca_{ij}$.
	\end{definition}
	
	\begin{proposition}{Properties of Matrix Operations}
		For matrices $A, B, C \in \R^{m \times n}$ and scalars $c, d \in \R$:
		\begin{enumerate}[label=(\roman*)]
			\item $A + B = B + A$ (commutativity)
			\item $(A + B) + C = A + (B + C)$ (associativity)
			\item $c(A + B) = cA + cB$ (distributivity)
			\item $(c + d)A = cA + dA$ (distributivity)
			\item $c(dA) = (cd)A$ (associativity of scalars)
		\end{enumerate}
	\end{proposition}
	
	\subsection{Matrix Multiplication}
	
	\begin{definition}{Matrix Multiplication}
		Let $A \in \R^{m \times n}$ and $B \in \R^{n \times p}$. The product $AB$ is the $m \times p$ matrix whose $(i,j)$-entry is:
		\[
		(AB)_{ij} = \sum_{k=1}^{n} a_{ik}b_{kj}
		\]
	\end{definition}
	
	\begin{remark}{Matrix Multiplication is Not Commutative}
		In general, $AB \neq BA$. In fact, $AB$ may be defined while $BA$ is not, or both may be defined but produce different results.
	\end{remark}
	
	\begin{example}{Simple Matrix Multiplication}
		Let $A = \begin{pmatrix} 1 & 2 \\ 3 & 4 \end{pmatrix}$ and $B = \begin{pmatrix} 5 & 6 \\ 7 & 8 \end{pmatrix}$.
		
		Then:
		\[
		AB = \begin{pmatrix} 1 \cdot 5 + 2 \cdot 7 & 1 \cdot 6 + 2 \cdot 8 \\ 3 \cdot 5 + 4 \cdot 7 & 3 \cdot 6 + 4 \cdot 8 \end{pmatrix} = \begin{pmatrix} 19 & 22 \\ 43 & 50 \end{pmatrix}
		\]
	\end{example}
	
	\section{Elementary Row Operations}
	
	Elementary row operations are the foundation of row reduction algorithms. These operations preserve the solution set of the system of linear equations represented by the matrix.
	
	\begin{definition}{Elementary Row Operations}
		The three elementary row operations on a matrix are:
		\begin{enumerate}[label=(\roman*)]
			\item \textbf{Row Scaling (Type I):} Multiply a row by a nonzero scalar $c \in \R$, $c \neq 0$. Notation: $R_i \to cR_i$.
			\item \textbf{Row Addition (Type II):} Add a multiple of one row to another row. Notation: $R_i \to R_i + cR_j$ where $i \neq j$ and $c \in \R$.
			\item \textbf{Row Swap (Type III):} Interchange two rows. Notation: $R_i \leftrightarrow R_j$ where $i \neq j$.
		\end{enumerate}
	\end{definition}
	
	\begin{theorem}{Elementary Row Operations Preserve Solution Sets}
		If a matrix $B$ is obtained from matrix $A$ by a sequence of elementary row operations, then the systems of linear equations represented by $A$ and $B$ have the same solution set.
	\end{theorem}
	
	\begin{example}{Elementary Row Operations}
		Consider the augmented matrix:
		\[
		\begin{pmatrix} 2 & 4 & | & 6 \\ 1 & 3 & | & 5 \end{pmatrix}
		\]
		
		Applying $R_1 \leftrightarrow R_2$:
		\[
		\begin{pmatrix} 1 & 3 & | & 5 \\ 2 & 4 & | & 6 \end{pmatrix}
		\]
		
		Applying $R_2 \to R_2 - 2R_1$:
		\[
		\begin{pmatrix} 1 & 3 & | & 5 \\ 0 & -2 & | & -4 \end{pmatrix}
		\]
	\end{example}
	
	\section{Row Reduction and Row Echelon Form}
	
	\begin{definition}{Row Echelon Form (REF)}
		A matrix is in row echelon form if it satisfies:
		\begin{enumerate}[label=(\roman*)]
			\item All zero rows (rows with all entries equal to zero) appear at the bottom of the matrix.
			\item The first nonzero entry (called the pivot) in each nonzero row is to the right of the pivot in the row above it.
			\item All entries below a pivot are zero.
		\end{enumerate}
	\end{definition}
	
	\begin{definition}{Reduced Row Echelon Form (RREF)}
		A matrix is in reduced row echelon form if it is in row echelon form and additionally satisfies:
		\begin{enumerate}[label=(\roman*), resume]
			\item Each pivot is equal to 1.
			\item All entries above and below each pivot are zero.
		\end{enumerate}
	\end{definition}
	
	\begin{remark}{Uniqueness of RREF}
		Every matrix has a unique reduced row echelon form. In contrast, the row echelon form of a matrix is not unique.
	\end{remark}
	
	\begin{example}{Row Echelon Form vs. RREF}
		The matrix
		\[
		\begin{pmatrix} 1 & 2 & 3 \\ 0 & 4 & 5 \\ 0 & 0 & 6 \end{pmatrix}
		\]
		is in row echelon form but not in reduced row echelon form.
		
		The matrix
		\[
		\begin{pmatrix} 1 & 0 & 0 \\ 0 & 1 & 0 \\ 0 & 0 & 1 \end{pmatrix}
		\]
		is in reduced row echelon form.
	\end{example}
	
	\section{Gaussian Elimination}
	
	Gaussian elimination is an algorithm for transforming a matrix into row echelon form using elementary row operations.
	
	\begin{definition}{Gaussian Elimination Algorithm}
		Given an $m \times n$ matrix $A$, Gaussian elimination proceeds as follows:
		\begin{enumerate}
			\item Identify the leftmost nonzero column. This is the first pivot column.
			\item If the top entry in the pivot column is zero, swap rows to place a nonzero entry at the top.
			\item Use row addition operations to make all entries below the pivot equal to zero.
			\item Repeat the process for the remaining submatrix (ignoring the current row and all rows above it), moving to the right.
			\item Continue until the matrix is in row echelon form.
		\end{enumerate}
	\end{definition}
	
	\begin{example}{Gaussian Elimination}
		Consider the augmented matrix:
		\[
		\left(\begin{array}{ccc|c}
			2 & 1 & -1 & 8 \\
			-3 & -1 & 2 & -11 \\
			-2 & 1 & 2 & -3
		\end{array}\right)
		\]
		
		\textbf{Step 1:} Use $R_1$ as the pivot row. Eliminate below:
		\begin{align*}
			R_2 &\to R_2 + \tfrac{3}{2}R_1 \\
			R_3 &\to R_3 + R_1
		\end{align*}
		\[
		\left(\begin{array}{ccc|c}
			2 & 1 & -1 & 8 \\
			0 & 1/2 & 1/2 & 1 \\
			0 & 2 & 1 & 5
		\end{array}\right)
		\]
		
		\textbf{Step 2:} Use $R_2$ as the pivot row. Eliminate below:
		\begin{align*}
			R_3 &\to R_3 - 4R_2
		\end{align*}
		\[
		\left(\begin{array}{ccc|c}
			2 & 1 & -1 & 8 \\
			0 & 1/2 & 1/2 & 1 \\
			0 & 0 & -1 & 1
		\end{array}\right)
		\]
		
		This matrix is now in row echelon form.
	\end{example}
	
	\begin{proposition}{Rank of a Matrix}
		The rank of a matrix $A$, denoted $\rank(A)$, is the number of nonzero rows in any row echelon form of $A$. Equivalently, it is the number of pivot columns.
	\end{proposition}
	
	\subsection{Forward Elimination}
	
	Forward elimination is the process of transforming a matrix into row echelon form. The algorithm works systematically from the top-left corner, moving right and down.
	
	\begin{theorem}{Forward Elimination Termination}
		Every matrix can be reduced to row echelon form by Gaussian elimination in a finite number of steps.
	\end{theorem}
	
	\section{Gauss-Jordan Elimination}
	
	Gauss-Jordan elimination is an extension of Gaussian elimination that transforms a matrix into reduced row echelon form.
	
	\begin{definition}{Gauss-Jordan Elimination Algorithm}
		Gauss-Jordan elimination extends Gaussian elimination by:
		\begin{enumerate}
			\item First, apply Gaussian elimination to obtain row echelon form.
			\item Then, perform back substitution using row operations to create zeros \textbf{above} each pivot as well.
			\item Scale each pivot row so that the pivot equals 1.
		\end{enumerate}
	\end{definition}
	
	\begin{example}{Gauss-Jordan Elimination}
		Continuing from the previous example, we have:
		\[
		\left(\begin{array}{ccc|c}
			2 & 1 & -1 & 8 \\
			0 & 1/2 & 1/2 & 1 \\
			0 & 0 & -1 & 1
		\end{array}\right)
		\]
		
		\textbf{Step 1:} Scale rows to have pivots equal to 1:
		\begin{align*}
			R_1 &\to \tfrac{1}{2}R_1 \\
			R_2 &\to 2R_2 \\
			R_3 &\to -R_3
		\end{align*}
		\[
		\left(\begin{array}{ccc|c}
			1 & 1/2 & -1/2 & 4 \\
			0 & 1 & 1 & 2 \\
			0 & 0 & 1 & -1
		\end{array}\right)
		\]
		
		\textbf{Step 2:} Eliminate above the pivots, working from bottom-right to top-left:
		\begin{align*}
			R_2 &\to R_2 - R_3 \\
			R_1 &\to R_1 + \tfrac{1}{2}R_3
		\end{align*}
		\[
		\left(\begin{array}{ccc|c}
			1 & 1/2 & 0 & 7/2 \\
			0 & 1 & 0 & 3 \\
			0 & 0 & 1 & -1
		\end{array}\right)
		\]
		
		\textbf{Step 3:} Eliminate above remaining pivots:
		\begin{align*}
			R_1 &\to R_1 - \tfrac{1}{2}R_2
		\end{align*}
		\[
		\left(\begin{array}{ccc|c}
			1 & 0 & 0 & 2 \\
			0 & 1 & 0 & 3 \\
			0 & 0 & 1 & -1
		\end{array}\right)
		\]
		
		This matrix is now in reduced row echelon form. The solution to the system is $(x, y, z) = (2, 3, -1)$.
	\end{example}
	
	\subsection{Back Substitution}
	
	Back substitution is the process of solving for the variables starting from the last equation and working backwards.
	
	\begin{definition}{Back Substitution}
		Once a system is in row echelon form (or reduced row echelon form), back substitution involves:
		\begin{enumerate}
			\item Identifying free variables (columns without pivots).
			\item Solving for the leading variables in terms of the free variables, starting from the bottom equation.
		\end{enumerate}
	\end{definition}
	
	\begin{example}{Back Substitution with Free Variables}
		Consider the system in RREF:
		\[
		\left(\begin{array}{ccccc|c}
			1 & 0 & 2 & 0 & 3 & 5 \\
			0 & 1 & -1 & 0 & 4 & 2 \\
			0 & 0 & 0 & 1 & -2 & 3 \\
			0 & 0 & 0 & 0 & 0 & 0
		\end{array}\right)
		\]
		
		The pivot columns are 1, 2, and 4. The free variables are $x_3$ and $x_5$.
		
		From the RREF form:
		\begin{align*}
			x_1 &= 5 - 2x_3 - 3x_5 \\
			x_2 &= 2 + x_3 - 4x_5 \\
			x_4 &= 3 + 2x_5
		\end{align*}
		
		Setting $x_3 = s$ and $x_5 = t$ (parameters), the general solution is:
		\[
		\begin{pmatrix} x_1 \\ x_2 \\ x_3 \\ x_4 \\ x_5 \end{pmatrix} = \begin{pmatrix} 5 \\ 2 \\ 0 \\ 3 \\ 0 \end{pmatrix} + s\begin{pmatrix} -2 \\ 1 \\ 1 \\ 0 \\ 0 \end{pmatrix} + t\begin{pmatrix} -3 \\ -4 \\ 0 \\ 2 \\ 1 \end{pmatrix}
		\]
	\end{example}
	
	\section{Applications and Special Cases}
	
	\subsection{Solving Linear Systems}
	
	\begin{theorem}{Solution Structure}
		For a system $A\vv{x} = \vv{b}$:
		\begin{enumerate}[label=(\roman*)]
			\item If $\rank(A) < \rank([A|\vv{b}])$, the system is inconsistent (no solution).
			\item If $\rank(A) = \rank([A|\vv{b}]) = n$ (number of columns), the system has a unique solution.
			\item If $\rank(A) = \rank([A|\vv{b}]) < n$, the system has infinitely many solutions.
		\end{enumerate}
	\end{theorem}
	
	\subsection{Finding the Inverse of a Matrix}
	
	\begin{definition}{Matrix Inverse}
		An $n \times n$ matrix $A$ is invertible if there exists an $n \times n$ matrix $A^{-1}$ such that:
		\[
		AA^{-1} = A^{-1}A = I_n
		\]
		where $I_n$ is the $n \times n$ identity matrix.
	\end{definition}
	
	\begin{theorem}{Gauss-Jordan Method for Finding Inverses}
		To find $A^{-1}$ for an invertible $n \times n$ matrix $A$:
		\begin{enumerate}
			\item Form the augmented matrix $[A|I_n]$.
			\item Apply Gauss-Jordan elimination to transform this into $[I_n|A^{-1}]$.
			\item The right half of the resulting matrix is $A^{-1}$.
		\end{enumerate}
	\end{theorem}
	
	\begin{example}{Finding a Matrix Inverse}
		Let $A = \begin{pmatrix} 1 & 2 \\ 3 & 4 \end{pmatrix}$. We form:
		\[
		\left[\begin{array}{cc|cc}
			1 & 2 & 1 & 0 \\
			3 & 4 & 0 & 1
		\end{array}\right]
		\]
		
		Apply $R_2 \to R_2 - 3R_1$:
		\[
		\left[\begin{array}{cc|cc}
			1 & 2 & 1 & 0 \\
			0 & -2 & -3 & 1
		\end{array}\right]
		\]
		
		Apply $R_2 \to -\frac{1}{2}R_2$:
		\[
		\left[\begin{array}{cc|cc}
			1 & 2 & 1 & 0 \\
			0 & 1 & 3/2 & -1/2
		\end{array}\right]
		\]
		
		Apply $R_1 \to R_1 - 2R_2$:
		\[
		\left[\begin{array}{cc|cc}
			1 & 0 & -2 & 1 \\
			0 & 1 & 3/2 & -1/2
		\end{array}\right]
		\]
		
		Therefore, $A^{-1} = \begin{pmatrix} -2 & 1 \\ 3/2 & -1/2 \end{pmatrix}$.
	\end{example}
	
	\subsection{Determining Linear Dependence}
	
	\begin{definition}{Linear Independence and Dependence}
		A set of vectors $\{\vv{v}_1, \vv{v}_2, \ldots, \vv{v}_k\}$ is linearly independent if the only solution to:
		\[
		c_1\vv{v}_1 + c_2\vv{v}_2 + \cdots + c_k\vv{v}_k = \vzero
		\]
		is $c_1 = c_2 = \cdots = c_k = 0$. Otherwise, the set is linearly dependent.
	\end{definition}
	
	\begin{proposition}{Testing Linear Independence via Row Reduction}
		Arrange the vectors as columns of a matrix and perform row reduction. The vectors are linearly independent if and only if each column is a pivot column in the row echelon form.
	\end{proposition}
	
	\begin{corollary}{Rank and Nullity}
		For an $m \times n$ matrix $A$:
		\[
		\rank(A) + \nullity(A) = n
		\]
		where $\nullity(A)$ is the dimension of the null space (kernel) of $A$.
	\end{corollary}
	
	\newpage
	
	\section{Summary and Key Takeaways}
	
	Row reduction is the cornerstone of computational linear algebra. Through elementary row operations, we can:
	
	\begin{enumerate}
		\item Solve systems of linear equations.
		\item Determine if a system is consistent or inconsistent.
		\item Find the rank of a matrix and identify pivot and free variables.
		\item Compute matrix inverses.
		\item Test linear independence of vectors.
		\item Find bases for fundamental subspaces.
	\end{enumerate}
	
	\textbf{Key Algorithms:}
	\begin{itemize}
		\item \textbf{Gaussian Elimination:} Transforms a matrix into row echelon form (REF).
		\item \textbf{Gauss-Jordan Elimination:} Transforms a matrix into reduced row echelon form (RREF).
		\item \textbf{Back Substitution:} Extracts the solution from row echelon form.
	\end{itemize}
	
	The reduced row echelon form is unique for each matrix and provides complete information about the solution set of the corresponding linear system.
	
\end{document}